\documentclass[11pt]{article}

\usepackage[utf8]{inputenc}
% microtype 의 폰트 확장은 scalable 폰트를 요구한다. 기본 CM 이 비트맵(Type3)으로
% 잡히는 배포판에서는 lmodern 없이는 "auto expansion is only possible with
% scalable fonts" 로 컴파일이 죽는다.
\usepackage{lmodern}
\usepackage[T1]{fontenc}
\usepackage[margin=1in]{geometry}
\usepackage{amsmath,amssymb}
\usepackage{graphicx}
\usepackage{longtable,booktabs}
\usepackage{microtype}
% svgnames 를 줘야 teal 이 정의된다. 없으면 \cite 하나마다
% "Undefined color `teal'" 이 뜬다 (본편 산출물에서 732회).
\usepackage[svgnames]{xcolor}
% hyperref goes last; it makes \cite clickable through to the bibliography.
\usepackage[colorlinks=true,linkcolor=blue,citecolor=teal,urlcolor=magenta]{hyperref}

% pandoc emits \tightlist but expects the preamble to define it.
\providecommand{\tightlist}{%
  \setlength{\itemsep}{0pt}\setlength{\parskip}{0pt}}

% If the compile times out on Overleaf, comment out \tableofcontents below first.
\title{Comprehensive Survey of Edge Computing: Architecture, Applications, and Challenges}
\author{Generated by AutoSurvey}
\date{}

\begin{document}
\maketitle
\tableofcontents
\newpage

\section{Introduction to Edge Computing}

\subsection{Definition and Motivation of Edge Computing}

Edge computing is a novel computing paradigm that extends cloud resources to the edge of the network, enabling data processing and analysis closer to where the data is generated \cite{ref1}. This approach is motivated by the need to overcome the limitations of traditional cloud computing, which can be plagued by high latency, bandwidth constraints, and security concerns. By pushing computing resources to the edge of the network, edge computing reduces latency and improves real-time processing capabilities, making it an attractive solution for applications that require fast data processing and analysis.

The proliferation of Internet of Things (IoT) devices has been a key driver of the emergence of edge computing \cite{ref2}. These devices generate vast amounts of data that need to be processed in real-time, and edge computing addresses this challenge by enabling data processing and analysis closer to the source. This reduces the amount of data that needs to be transmitted to the cloud, thereby minimizing the risk of data breaches and cyber attacks \cite{ref3}. Additionally, edge computing enables real-time analytics and decision-making, which is critical for applications such as autonomous vehicles, smart cities, and industrial automation \cite{ref4}.

The concept of edge computing has gained significant traction in recent years due to advances in technologies such as 5G networks, artificial intelligence, and the IoT \cite{ref5}. As a result, edge computing is now recognized as a critical component of the digital transformation journey, enabling organizations to unlock new revenue streams, improve customer experiences, and drive business innovation. Edge computing has several key characteristics that distinguish it from traditional cloud computing, including low latency, real-time processing, and enhanced security, as well as the ability to operate in resource-constrained environments \cite{ref6}.

In terms of its capabilities, edge computing enables new use cases and applications, such as augmented reality, virtual reality, and machine learning, which require real-time processing and analysis of data \cite{ref7}. By extending cloud resources to the edge of the network, edge computing provides a powerful platform for data processing and analysis, enabling organizations to extract insights and make informed decisions in real-time. In summary, edge computing is a novel computing paradigm that extends cloud resources to the edge of the network, enabling data processing and analysis closer to where the data is generated \cite{ref8}.

\subsection{Key Concepts and Definitions in Edge Computing}

The field of edge computing has been extensively explored in various research papers, each contributing to a deeper understanding of its key concepts, definitions, and applications. \cite{ref1} provides a comprehensive overview of the current initiatives and a roadmap for sustainable edge computing development, highlighting the importance of edge computing in reducing latency and improving real-time processing. This is further supported by \cite{ref8}, which conducts a systematic mapping study to identify the main topics, architectures, techniques, and applications in the field of edge computing, providing a thorough understanding of the field's current state.

One of the fundamental concepts in edge computing is the idea of pushing computing resources to the edge of the network, closer to the devices and data sources. This approach is critical in addressing the security threats and vulnerabilities emerging from the edge network architecture, as discussed in \cite{ref3}, which highlights the need for a holistic approach to analyze edge network security posture. Furthermore, the integration of artificial intelligence (AI) with edge computing, referred to as Edge AI, is presented in \cite{ref9}, which provides a comprehensive discussion on AI methods and capabilities in the context of edge computing.

The literature surrounding edge intelligence is comprehensively surveyed in \cite{ref10}, which identifies four fundamental components of edge intelligence: edge caching, edge training, edge inference, and edge offloading. Additionally, \cite{ref11} provides a comprehensive overview of the current state of research in edge-to-cloud computing, while \cite{ref12} examines the history of benchmarking from tightly coupled to loosely coupled systems. The considerations for edge deployments in B5G/6G networks are presented in \cite{ref13}.

The exploration of edge computing concepts and definitions is further enriched by various studies, including \cite{ref14}, which provides a taxonomy for edge computing, and \cite{ref15}, which discusses the security issues and challenges in edge computing. Moreover, recent advances in binary neural networks for edge computing are reviewed in \cite{ref16}, and the emerging research area of edge intelligence is overviewed in \cite{ref17}.

The role of edge computing in realizing the vision of smart cities is highlighted in \cite{ref18}, and the conditions under which the edge is feasible are investigated in \cite{ref19}. The major constraints in operationalizing Deep Edge Intelligence (DEI) are overviewed in \cite{ref20}, and recent access control measures are reviewed in \cite{ref21}. Furthermore, the characteristics, architecture, technologies, and solutions of trustworthy Edge Intelligence (EI) are comprehensively summarized in \cite{ref22}, and an Open Framework for Edge Intelligence (OpenEI) is presented in \cite{ref23}.

In conclusion, the key concepts and definitions in edge computing have been extensively explored in various papers, highlighting the importance of edge computing in reducing latency, improving real-time processing, and enabling novel applications and services. These studies collectively contribute to a deeper understanding of edge computing and its potential to transform various aspects of modern computing and communication systems.

\section{Edge Computing Architecture and Technologies}

\subsection{Edge Computing Architecture}

The architecture of edge computing is a complex system that enables the distribution of computing resources and services at the edge of the network, closer to the source of the data. This paradigm shift from traditional cloud computing to edge computing is driven by the need for lower latency, improved real-time processing, and enhanced security \cite{ref1}. The edge computing architecture consists of various components, including edge devices, edge servers, and communication protocols, which work together to provide a robust and scalable platform for deploying applications and services.

Edge devices are the endpoints of the edge computing system, which can be anything from smartphones and laptops to IoT devices and sensors \cite{ref2}. These devices generate data that needs to be processed and analyzed in real-time, and they often have limited computing resources and storage capacity. Edge servers, on the other hand, are more powerful computing devices that are located at the edge of the network, closer to the edge devices \cite{ref24}. They provide the necessary computing resources and storage capacity to process and analyze the data generated by the edge devices, and are often used in conjunction with enabling technologies such as containerization, virtualization, and fog computing.

The communication protocols used in edge computing are designed to enable efficient and reliable communication between the edge devices and edge servers \cite{ref25}. These protocols include wireless communication protocols such as Wi-Fi, Bluetooth, and cellular networks, as well as wired communication protocols such as Ethernet. The choice of communication protocol depends on the specific use case and the requirements of the application, and is critical to ensuring the smooth operation of the edge computing system.

One of the key challenges in edge computing is the management of the edge devices and edge servers \cite{ref26}. This includes tasks such as device discovery, device management, and resource allocation. Device discovery involves identifying and registering the edge devices on the network, while device management involves monitoring and controlling the edge devices to ensure they are functioning correctly. Resource allocation involves allocating the necessary computing resources and storage capacity to the edge devices and edge servers, and is critical to ensuring the efficient operation of the edge computing system.

Another important aspect of edge computing is security \cite{ref27}. Edge computing systems are vulnerable to various types of attacks, including data breaches, denial-of-service attacks, and malware attacks. To mitigate these risks, edge computing systems use various security measures such as encryption, authentication, and access control. Encryption involves encrypting the data transmitted between the edge devices and edge servers, while authentication involves verifying the identity of the edge devices and edge servers. Access control involves controlling who has access to the edge devices and edge servers, and is critical to ensuring the security and integrity of the edge computing system.

In addition to security, edge computing also raises concerns about privacy \cite{ref8}. Edge computing systems often collect and process sensitive data, including personal data and confidential business information. To protect this data, edge computing systems use various privacy measures such as data anonymization, data aggregation, and data encryption. Data anonymization involves removing personally identifiable information from the data, while data aggregation involves combining data from multiple sources to prevent individual identification. Data encryption involves encrypting the data to prevent unauthorized access, and is critical to ensuring the privacy and confidentiality of the data.

In conclusion, the architecture of edge computing is a complex system that consists of edge devices, edge servers, and communication protocols, which work together to provide a robust and scalable platform for deploying applications and services. The management of edge devices and edge servers, security, and privacy are critical aspects of edge computing that need to be carefully considered, and are essential to ensuring the smooth operation and integrity of the edge computing system \cite{ref28}. By understanding the architecture of edge computing and the various technologies that enable it, we can better appreciate the benefits and challenges of this emerging technology and develop effective solutions to address the challenges.

\subsection{Enabling Technologies for Edge Computing}

Edge computing is enabled by a variety of technologies that facilitate the deployment, management, and execution of applications at the edge of the network. As discussed in the previous section, the architecture of edge computing consists of edge devices, edge servers, and communication protocols, which work together to provide a robust and scalable platform for deploying applications and services. To further enhance the capabilities of edge computing, several enabling technologies have been developed, including containerization \cite{ref29}, which allows for the packaging of applications and their dependencies into a single container, providing a high degree of flexibility and scalability.

In addition to containerization, virtualization \cite{ref30} is another important technology for edge computing, enabling the creation of virtual machines (VMs) that can run on edge devices, providing a layer of abstraction between the application and the underlying hardware. This allows for the deployment of multiple VMs on a single edge device, each running a different application or service. Fog computing \cite{ref31} is also a key enabling technology for edge computing, extending the cloud computing paradigm to the edge of the network, providing a hierarchical architecture for computing, storage, and networking.

Other technologies, such as cloudlets \cite{ref32}, which are small-scale cloud data centers deployed at the edge of the network, and mobile edge computing (MEC) \cite{ref24}, which provides a platform for deploying applications and services at the edge of the network, using the resources of mobile devices and edge servers, also play a crucial role in enabling edge computing. Furthermore, software-defined networking (SDN) and network function virtualization (NFV) \cite{ref33} enable the creation of virtual networks and network functions, which can be deployed at the edge of the network, providing a high degree of flexibility and scalability.

The integration of these technologies enables the creation of a robust and scalable edge computing environment, which can support a wide range of applications and services \cite{ref34}. This environment can be used to deploy applications and services at the edge of the network, reducing latency and improving real-time processing capabilities. The use of containerization, virtualization, fog computing, cloudlets, and MEC enables the creation of a flexible and scalable edge computing environment, which can support a wide range of use cases and applications, and provide a high degree of flexibility and scalability, making it ideal for edge computing environments \cite{ref1}. As we will discuss in the next section, the applications and use cases of edge computing are diverse and continue to grow, driven by the increasing demand for low-latency and real-time processing capabilities.

\section{Applications and Use Cases of Edge Computing}

\subsection{Edge Computing Applications in IoT and Smart Cities}

Edge computing has numerous applications in the Internet of Things (IoT) and smart cities, where it can significantly enhance the efficiency, security, and scalability of various services. In IoT, edge computing can be used to process data generated by sensors and devices in real-time, reducing latency and improving the overall performance of the system \cite{ref2}. For instance, in industrial automation, edge computing can be used to monitor and control equipment, predict maintenance needs, and optimize production processes \cite{ref35}. This enables businesses to respond quickly to changes in the production process, reduce downtime, and improve overall productivity.

In smart cities, edge computing can be used to manage and analyze data from various sources, such as traffic cameras, sensors, and social media, to improve public safety, transportation, and energy management \cite{ref18}. For example, edge computing can be used to analyze video feeds from traffic cameras to detect accidents, traffic congestion, and other incidents, and alert authorities in real-time \cite{ref36}. This enables city officials to respond quickly to emergencies, reduce traffic congestion, and improve the overall quality of life for citizens.

Edge computing can also be used to support various smart city applications, such as smart parking, smart lighting, and smart waste management \cite{ref37}. For instance, edge computing can be used to analyze data from sensors and cameras to detect available parking spots, and guide drivers to them, reducing traffic congestion and pollution \cite{ref38}. Additionally, edge computing can be used to optimize energy consumption in smart buildings, reducing energy waste and improving the overall sustainability of the city.

Furthermore, edge computing can be used to enhance the security and privacy of IoT devices and smart city infrastructure \cite{ref39}. For example, edge computing can be used to detect and respond to cyber threats in real-time, reducing the risk of data breaches and other security incidents \cite{ref40}. This enables cities to protect sensitive data, prevent cyber attacks, and maintain the trust of citizens.

In addition, edge computing can be used to support the development of autonomous vehicles, which require low-latency and high-bandwidth communications to operate safely and efficiently \cite{ref41}. Edge computing can be used to process data from sensors and cameras, and make decisions in real-time, reducing the risk of accidents and improving the overall safety of the vehicle \cite{ref42}. This is particularly important in smart cities, where autonomous vehicles can improve traffic flow, reduce congestion, and enhance the overall quality of life for citizens.

Overall, edge computing has the potential to transform the way IoT and smart city applications are designed, deployed, and managed, enabling new use cases, improving efficiency, and reducing costs \cite{ref1}. As the number of IoT devices and smart city applications continues to grow, edge computing is likely to play an increasingly important role in supporting these applications, and enabling new innovations and services \cite{ref43}. By leveraging edge computing, businesses and cities can unlock new opportunities, improve operational efficiency, and create a more sustainable and connected future.

\subsection{Edge Computing Use Cases in Healthcare, Gaming, and Autonomous Vehicles}

Edge computing has numerous use cases in various industries, including healthcare, gaming, and autonomous vehicles, in addition to its applications in IoT and smart cities. Building on its potential to enhance efficiency, security, and scalability in IoT and smart city applications, edge computing can be used to analyze medical images, such as X-rays and MRIs, in real-time, enabling doctors to make quick and accurate diagnoses \cite{ref44}. For instance, edge devices can be used to detect abnormalities in medical images, such as tumors or fractures, and alert doctors to take immediate action. Additionally, edge computing can be used to monitor patients' vital signs, such as heart rate and blood pressure, in real-time, enabling healthcare professionals to respond quickly to any changes in their condition \cite{ref45}.

In the gaming industry, edge computing can be used to reduce latency and improve the overall gaming experience \cite{ref46}. For example, edge computing can be used to render graphics in real-time, enabling gamers to enjoy a more immersive and responsive gaming experience. Moreover, edge computing can be used to analyze gamer behavior and provide personalized recommendations, such as suggesting new games or offering tips to improve their gaming skills \cite{ref2}. This application of edge computing in gaming is a natural extension of its use in IoT and smart cities, where real-time data analysis and processing are critical to improving efficiency and decision-making.

Autonomous vehicles are another area where edge computing can be applied \cite{ref47}. Edge computing can be used to analyze sensor data from cameras, lidar, and radar in real-time, enabling autonomous vehicles to detect and respond to their surroundings quickly and accurately. For instance, edge computing can be used to detect pedestrians, traffic lights, and other obstacles, and adjust the vehicle's trajectory accordingly. Additionally, edge computing can be used to optimize route planning and traffic flow, reducing congestion and improving overall traffic efficiency \cite{ref48}. The use of edge computing in autonomous vehicles is closely related to its application in smart cities, where real-time data analysis and processing are critical to improving traffic flow and reducing congestion.

The use of edge computing in autonomous vehicles can also improve safety \cite{ref49}. For example, edge computing can be used to detect potential accidents, such as a pedestrian stepping into the road, and alert the driver or take control of the vehicle to avoid the accident. Moreover, edge computing can be used to analyze data from various sources, such as traffic cameras and sensors, to predict and prevent accidents. This application of edge computing in autonomous vehicles highlights its potential to transform various industries by enabling real-time data analysis, reducing latency, and improving overall efficiency \cite{ref44}.

In addition to these use cases, edge computing can also be used in other areas, such as smart cities and industrial automation \cite{ref50}. For instance, edge computing can be used to analyze data from sensors and cameras in smart cities, enabling city officials to monitor and manage traffic flow, energy usage, and waste management more efficiently. In industrial automation, edge computing can be used to analyze data from sensors and machines, enabling manufacturers to optimize production processes, predict maintenance needs, and improve overall efficiency. As the amount of data generated by devices and sensors continues to grow, edge computing will play an increasingly important role in analyzing and processing this data, enabling businesses and organizations to make better decisions and improve their operations, and paving the way for further innovations and applications in various industries.

\section{Challenges and Open Research Issues in Edge Computing}

\subsection{Security and Privacy Challenges in Edge Computing}

Security and privacy are significant concerns in edge computing, as the paradigm involves processing and storing sensitive data at the edge of the network, closer to the source of the data \cite{ref3}. The distributed nature of edge computing, with multiple edge devices and servers, further complicates the security landscape, making it challenging to ensure the confidentiality, integrity, and availability of data \cite{ref51}. As edge computing continues to grow and expand into various applications, including IoT, smart cities, and healthcare, the need for robust security and privacy measures becomes increasingly important.

One of the primary security challenges in edge computing is the risk of data leakage and unauthorized access to sensitive information \cite{ref52}. This risk is particularly pronounced in applications that involve sensitive personal data, such as healthcare and finance, where the consequences of a data breach can be severe \cite{ref53}. Furthermore, the use of edge computing in IoT applications, which often involve a large number of devices, increases the attack surface and makes it more challenging to ensure the security of the system \cite{ref54}.

In addition to security challenges, edge computing also raises significant privacy concerns \cite{ref55}. The processing and storage of sensitive data at the edge of the network, closer to the source of the data, can compromise the privacy of individuals and organizations. Moreover, the use of edge computing in applications such as surveillance and monitoring can raise concerns about data privacy and surveillance \cite{ref56}. To mitigate these concerns, it is essential to develop and implement effective security and privacy solutions that can ensure the confidentiality, integrity, and availability of data, while also protecting the privacy of individuals and organizations.

Various solutions have been proposed to address the security and privacy challenges in edge computing, including the use of blockchain technology \cite{ref57}, secure multi-party computation \cite{ref58}, and homomorphic encryption \cite{ref59}. These solutions aim to provide a robust and secure framework for edge computing, enabling the widespread adoption of the paradigm in various applications. By leveraging these technologies, it is possible to ensure the security and privacy of edge computing systems and protect the sensitive data processed and stored at the edge of the network.

In conclusion, security and privacy are critical components of edge computing, and addressing these challenges is essential to ensuring the success of the paradigm. As edge computing continues to evolve and expand into new applications, the need for robust security and privacy measures will only continue to grow. By developing and implementing effective security and privacy solutions, we can unlock the full potential of edge computing and enable the creation of secure, private, and efficient edge-based systems \cite{ref60}. This will be crucial in addressing the resource management and optimization challenges in edge computing, which will be discussed in the following section.

\subsection{Resource Management and Optimization Challenges in Edge Computing}

Resource management and optimization are crucial challenges in edge computing, as they directly impact the performance, efficiency, and scalability of edge-based systems. As discussed in the previous section, security and privacy are critical components of edge computing, and addressing these challenges is essential to ensuring the success of the paradigm. Similarly, effective resource management and optimization are necessary to unlock the full potential of edge computing and enable the creation of efficient, scalable, and sustainable edge-based systems. Edge computing involves deploying computing resources at the edge of the network, closer to the users and devices, to reduce latency and improve real-time processing. However, this paradigm also introduces new resource management challenges, such as managing heterogeneous resources, handling dynamic workloads, and optimizing resource allocation \cite{ref61}.

One of the key resource management challenges in edge computing is the allocation of computing resources, such as CPU, memory, and storage, to various applications and services \cite{ref62}. This challenge is exacerbated by the fact that edge devices are often resource-constrained, and the availability of resources can vary significantly depending on the location, time, and other factors \cite{ref63}. To address this challenge, researchers have proposed various resource allocation strategies, including dynamic resource allocation, resource provisioning, and workload scheduling \cite{ref64}. These strategies aim to optimize resource utilization, reduce latency, and improve the overall performance of edge-based systems.

Another significant challenge in edge computing is the management of network resources, such as bandwidth, latency, and packet loss \cite{ref65}. The allocation of network resources is critical to ensure that edge-based applications and services can communicate efficiently and effectively with each other and with the cloud \cite{ref66}. Researchers have proposed various network resource management strategies, including network slicing, traffic engineering, and quality of service (QoS) management \cite{ref67}. These strategies aim to optimize network resource utilization, reduce congestion, and improve the overall performance of edge-based systems.

In addition to computing and network resource management, edge computing also requires effective management of energy resources \cite{ref68}. Edge devices are often powered by batteries or other energy-harvesting technologies, and minimizing energy consumption is critical to prolong the lifetime of these devices \cite{ref69}. Researchers have proposed various energy management strategies, including energy-aware resource allocation, energy harvesting, and power management \cite{ref70}. These strategies aim to optimize energy utilization, reduce energy consumption, and improve the overall sustainability of edge-based systems.

To address the resource management and optimization challenges in edge computing, researchers have proposed various frameworks, architectures, and algorithms \cite{ref71}. These include edge node placement, resource allocation, and workload scheduling \cite{ref72}. Machine learning and artificial intelligence (AI) techniques have also been applied to edge computing resource management, including predictive modeling, reinforcement learning, and deep learning \cite{ref73}. These techniques aim to optimize resource utilization, improve system performance, and enable more efficient and scalable edge-based systems.

In conclusion, resource management and optimization are critical challenges in edge computing, and addressing these challenges requires a comprehensive understanding of the complex interactions between computing, network, and energy resources. By developing and implementing effective resource management and optimization strategies, we can unlock the full potential of edge computing and enable the creation of efficient, scalable, and sustainable edge-based systems. As edge computing continues to evolve and expand into new applications and use cases, the need for robust resource management and optimization solutions will only continue to grow, and further research is needed to explore new approaches and techniques for addressing these challenges \cite{ref74}. The development of these solutions will be crucial in enabling the widespread adoption of edge computing and unlocking its full potential to transform various industries and applications \cite{ref75}.

\section{Future Directions and Conclusion}

\subsection{Future Directions of Edge Computing}

The future of edge computing holds immense promise, with its potential to revolutionize various aspects of our lives, from industries to personal experiences. As edge computing continues to evolve, its integration with emerging technologies is expected to play a crucial role in shaping its future directions. One of the key areas of focus will be the convergence of edge computing with artificial intelligence (AI) and machine learning (ML), enabling the development of more intelligent and autonomous systems \cite{ref28}. This integration will facilitate the creation of edge AI applications that can process vast amounts of data in real-time, making them ideal for applications such as smart cities, healthcare, and transportation.

The combination of edge computing and other emerging technologies, such as the Internet of Things (IoT), will also be a significant direction for future research and development \cite{ref2}. By processing data closer to where it is generated, edge computing can reduce latency and improve the overall performance of IoT applications. Furthermore, the use of edge computing in IoT will also enable the development of more secure and reliable systems, as data will not need to be transmitted to the cloud or centralized servers for processing. The emergence of 5G and 6G networks will provide the necessary infrastructure for the widespread adoption of edge computing, enabling the creation of more responsive and efficient edge computing systems \cite{ref76}.

In addition to its integration with AI, ML, and IoT, edge computing will also be shaped by its convergence with other emerging technologies, such as blockchain and quantum computing \cite{ref77}. The use of blockchain in edge computing will enable the creation of more secure and transparent systems, as data can be stored and transmitted in a secure and decentralized manner. As edge computing continues to evolve, it is expected to have a significant impact on various industries, including healthcare, manufacturing, and transportation \cite{ref78}. The use of edge computing in these industries will enable the creation of more efficient and effective systems, as data can be processed and analyzed in real-time, enabling faster diagnosis, treatment, and decision-making.

The development of new edge computing architectures and platforms, such as microservice-based edge platforms and serverless edge computing, will also play a crucial role in shaping the future of edge computing \cite{ref79}. These new architectures and platforms will enable the creation of more efficient and effective edge computing systems, making them ideal for various applications. As we move forward, it is essential to continue researching and developing edge computing technologies, to unlock its full potential and enable the creation of more efficient, effective, and secure systems. By doing so, we can ensure that edge computing reaches its full potential and has a positive impact on various aspects of our lives, from industries to personal experiences.

\subsection{Conclusion and Open Research Issues}

In conclusion, this survey has provided a comprehensive overview of the current state of edge computing, including its architecture, applications, and challenges. The key findings of this survey highlight the importance of edge computing in reducing latency, improving real-time processing, and enhancing security \cite{ref3}. As discussed in the previous sections, the integration of edge computing with emerging technologies such as AI, ML, and IoT is expected to play a crucial role in shaping its future directions.

One of the major challenges in edge computing is the lack of standardization, which hinders the widespread adoption of edge computing technologies \cite{ref1}. Furthermore, the survey highlights the need for more research on edge computing applications, such as IoT, smart cities, industrial automation, healthcare, transportation, gaming, and autonomous vehicles, which will be enabled by the convergence of edge computing with other emerging technologies.

The emergence of new technologies, such as 5G and 6G networks, is expected to further accelerate the adoption of edge computing \cite{ref13}. However, this will also introduce new challenges, such as the need for more advanced security measures and more efficient resource management \cite{ref80}. To address these challenges, it is essential to continue researching and developing edge computing technologies, including the development of new edge computing architectures and platforms.

In addition, the survey highlights the importance of edge intelligence, which refers to the integration of edge computing and artificial intelligence \cite{ref81}. Edge intelligence has the potential to enable a wide range of applications, such as real-time data processing, autonomous vehicles, and smart cities \cite{ref28}. The development of edge intelligence will be critical in unlocking the full potential of edge computing and enabling the creation of more efficient, effective, and secure systems.

The survey also identifies several open research issues and challenges in edge computing, such as the need for more efficient resource management, more advanced security measures, and more efficient data processing \cite{ref6}. Furthermore, the survey highlights the need for more research on edge computing applications, such as IoT, smart cities, industrial automation, healthcare, transportation, gaming, and autonomous vehicles \cite{ref21}. These open research issues and challenges will be addressed through continued research and development, including the exploration of new edge computing architectures and platforms.

In terms of future research directions, the survey suggests that more research is needed on edge computing applications, such as IoT, smart cities, industrial automation, healthcare, transportation, gaming, and autonomous vehicles \cite{ref4}. Additionally, more research is needed on edge intelligence, including the development of more advanced AI models and more efficient data processing \cite{ref82}. The future of edge computing holds immense promise, and it is essential to continue researching and developing edge computing technologies to unlock its full potential.

Finally, the survey highlights the need for more collaboration between industry and academia to address the challenges and open research issues in edge computing \cite{ref83}. By working together, we can ensure that edge computing reaches its full potential and has a positive impact on various aspects of our lives, from industries to personal experiences. Overall, this survey provides a comprehensive overview of the current state of edge computing and identifies several open research issues and challenges that need to be addressed, highlighting the importance of edge computing in reducing latency, improving real-time processing, and enhancing security, and suggesting several future research directions.


\begin{thebibliography}{83}
\addcontentsline{toc}{section}{References}
\bibitem{ref1} Edge Computing: A Comprehensive Survey of Current Initiatives and a Roadmap for a Sustainable Edge Computing Development. arXiv:1912.08530, 2019. \url{https://arxiv.org/abs/1912.08530}
\bibitem{ref2} Edge Computing for IoT. arXiv:2402.13056, 2024. \url{https://arxiv.org/abs/2402.13056}
\bibitem{ref3} Edge Security: Challenges and Issues. arXiv:2206.07164, 2022. \url{https://arxiv.org/abs/2206.07164}
\bibitem{ref4} Towards Enabling Novel Edge-Enabled Applications. arXiv:1805.06989, 2018. \url{https://arxiv.org/abs/1805.06989}
\bibitem{ref5} Revisiting the Arguments for Edge Computing Research. arXiv:2106.12224, 2021. \url{https://arxiv.org/abs/2106.12224}
\bibitem{ref6} Dependability in Edge Computing. arXiv:1710.11222, 2017. \url{https://arxiv.org/abs/1710.11222}
\bibitem{ref7} Deep Learning at the Edge. arXiv:1910.10231, 2019. \url{https://arxiv.org/abs/1910.10231}
\bibitem{ref8} Edge Computing: A Systematic Mapping Study. arXiv:2102.02720, 2021. \url{https://arxiv.org/abs/2102.02720}
\bibitem{ref9} Roadmap for Edge AI: A Dagstuhl Perspective. arXiv:2112.00616, 2021. \url{https://arxiv.org/abs/2112.00616}
\bibitem{ref10} Edge Intelligence: Architectures, Challenges, and Applications. arXiv:2003.12172, 2020. \url{https://arxiv.org/abs/2003.12172}
\bibitem{ref11} Edge to Cloud Tools: A Multivocal Literature Review. arXiv:2305.17464, 2023. \url{https://arxiv.org/abs/2305.17464}
\bibitem{ref12} A Survey on Edge Benchmarking. arXiv:2004.11725, 2020. \url{https://arxiv.org/abs/2004.11725}
\bibitem{ref13} Edge Computing in IoT: A 6G Perspective. arXiv:2111.08943, 2021. \url{https://arxiv.org/abs/2111.08943}
\bibitem{ref14} A Taxonomy for Management and Optimization of Multiple Resources in Edge Computing. arXiv:1801.05610, 2018. \url{https://arxiv.org/abs/1801.05610}
\bibitem{ref15} Edge Computing in Transportation: Security Issues and Challenges. arXiv:2012.11206, 2020. \url{https://arxiv.org/abs/2012.11206}
\bibitem{ref16} A Review of Recent Advances of Binary Neural Networks for Edge Computing. arXiv:2011.14824, 2020. \url{https://arxiv.org/abs/2011.14824}
\bibitem{ref17} Pushing AI to Wireless Network Edge: An Overview on Integrated Sensing, Communication, and Computation towards 6G. arXiv:2211.02574, 2022. \url{https://arxiv.org/abs/2211.02574}
\bibitem{ref18} Edge-Computing-Enabled Smart Cities: A Comprehensive Survey. arXiv:1909.08747, 2019. \url{https://arxiv.org/abs/1909.08747}
\bibitem{ref19} To Edge or Not to Edge?. arXiv:1803.08448, 2018. \url{https://arxiv.org/abs/1803.08448}
\bibitem{ref20} State-of-the-art Techniques in Deep Edge Intelligence. arXiv:2008.00824, 2020. \url{https://arxiv.org/abs/2008.00824}
\bibitem{ref21} Edge Computing for IoT: Novel Insights from a Comparative Analysis of Access Control Models. arXiv:2405.07685, 2024. \url{https://arxiv.org/abs/2405.07685}
\bibitem{ref22} A Survey on Trustworthy Edge Intelligence: From Security and Reliability To Transparency and Sustainability. arXiv:2310.17944, 2023. \url{https://arxiv.org/abs/2310.17944}
\bibitem{ref23} OpenEI: An Open Framework for Edge Intelligence. arXiv:1906.01864, 2019. \url{https://arxiv.org/abs/1906.01864}
\bibitem{ref24} Mobile Edge Computing. arXiv:2404.17944, 2024. \url{https://arxiv.org/abs/2404.17944}
\bibitem{ref25} Wireless Edge Computing With Latency and Reliability Guarantees. arXiv:1905.05316, 2019. \url{https://arxiv.org/abs/1905.05316}
\bibitem{ref26} Resource Management in Edge and Fog Computing using FogBus2 Framework. arXiv:2108.00591, 2021. \url{https://arxiv.org/abs/2108.00591}
\bibitem{ref27} The Security and Privacy of Mobile Edge Computing: An Artificial Intelligence Perspective. arXiv:2401.01589, 2024. \url{https://arxiv.org/abs/2401.01589}
\bibitem{ref28} Edge AI: A Taxonomy, Systematic Review and Future Directions. arXiv:2407.04053, 2024. \url{https://arxiv.org/abs/2407.04053}
\bibitem{ref29} Edge System Design Using Containers and Unikernels for IoT Applications. arXiv:2412.03032, 2024. \url{https://arxiv.org/abs/2412.03032}
\bibitem{ref30} Fog computing state of the art: concept and classification of platforms to support distributed computing systems. arXiv:2106.11726, 2021. \url{https://arxiv.org/abs/2106.11726}
\bibitem{ref31} Fog Computing: Focusing on Mobile Users at the Edge. arXiv:1502.01815, 2015. \url{https://arxiv.org/abs/1502.01815}
\bibitem{ref32} Edge, Fog, and Cloud Computing : An Overview on Challenges and Applications. arXiv:2211.01863, 2022. \url{https://arxiv.org/abs/2211.01863}
\bibitem{ref33} A View of Fog Computing from Networking Perspective. arXiv:1602.01509, 2016. \url{https://arxiv.org/abs/1602.01509}
\bibitem{ref34} Integration of FogBus2 Framework with Container Orchestration Tools in Cloud and Edge Computing Environments. arXiv:2112.02267, 2021. \url{https://arxiv.org/abs/2112.02267}
\bibitem{ref35} Emerging Edge Computing Technologies for Distributed IoT Systems. arXiv:1811.11268, 2018. \url{https://arxiv.org/abs/1811.11268}
\bibitem{ref36} FOCAN: A Fog-supported smart city network architecture for management of applications in the Internet of Everything environments. arXiv:1710.01801, 2017. \url{https://arxiv.org/abs/1710.01801}
\bibitem{ref37} A Commodity SBC-Edge Cluster for Smart Cities. arXiv:1902.06661, 2019. \url{https://arxiv.org/abs/1902.06661}
\bibitem{ref38} SAFEMYRIDES: Application of Decentralized Control Edge-Computing to Ridesharing Monitoring Services. arXiv:2301.00888, 2023. \url{https://arxiv.org/abs/2301.00888}
\bibitem{ref39} Secure solutions for Smart City Command Control Centre using AIOT. arXiv:2108.00003, 2021. \url{https://arxiv.org/abs/2108.00003}
\bibitem{ref40} Mobile Privacy-Preserving Crowdsourced Data Collection in the Smart City. arXiv:1607.02805, 2016. \url{https://arxiv.org/abs/1607.02805}
\bibitem{ref41} Deep Reinforcement Learning for Adaptive Network Slicing in 5G for Intelligent Vehicular Systems and Smart Cities. arXiv:2010.09916, 2020. \url{https://arxiv.org/abs/2010.09916}
\bibitem{ref42} A Novel Design for Advanced 5G Deployment Environments with Virtualized Resources at Vehicular and MEC Nodes. arXiv:2303.15836, 2023. \url{https://arxiv.org/abs/2303.15836}
\bibitem{ref43} Quantum-Edge Cloud Computing: A Future Paradigm for IoT Applications. arXiv:2405.04824, 2024. \url{https://arxiv.org/abs/2405.04824}
\bibitem{ref44} Edge Computing for Smart Health: Context-Aware Approaches, Opportunities, and Challenges. arXiv:2004.07311, 2020. \url{https://arxiv.org/abs/2004.07311}
\bibitem{ref45} Role of Edge Device and Cloud Machine Learning in Point-of-Care Solutions Using Imaging Diagnostics for Population Screening. arXiv:2006.13808, 2020. \url{https://arxiv.org/abs/2006.13808}
\bibitem{ref46} AI on the Edge: Rethinking AI-based IoT Applications Using Specialized Edge Architectures. arXiv:2003.12488, 2020. \url{https://arxiv.org/abs/2003.12488}
\bibitem{ref47} Digital Twins for Autonomous Driving: A Comprehensive Implementation and Demonstration. arXiv:2401.08653, 2024. \url{https://arxiv.org/abs/2401.08653}
\bibitem{ref48} Vehicular Cloud Computing: A cost-effective alternative to Edge Computing in 5G networks. arXiv:2507.15670, 2025. \url{https://arxiv.org/abs/2507.15670}
\bibitem{ref49} Preventing and Controlling Epidemics through Blockchain-Assisted AI-Enabled Networks. arXiv:2104.13224, 2021. \url{https://arxiv.org/abs/2104.13224}
\bibitem{ref50} Industrial Edge-based Cyber-Physical Systems -- Application Needs and Concerns for Realization. arXiv:2201.00242, 2022. \url{https://arxiv.org/abs/2201.00242}
\bibitem{ref51} Security in Mobile Edge Caching with Reinforcement Learning. arXiv:1801.05915, 2018. \url{https://arxiv.org/abs/1801.05915}
\bibitem{ref52} Privacy-Preserving Detection Method for Transmission Line Based on Edge Collaboration. arXiv:2308.08761, 2023. \url{https://arxiv.org/abs/2308.08761}
\bibitem{ref53} Blockchain Powered Edge Intelligence for U-Healthcare in Privacy Critical and Time Sensitive Environment. arXiv:2506.02038, 2025. \url{https://arxiv.org/abs/2506.02038}
\bibitem{ref54} Security and Privacy for Mobile Edge Caching: Challenges and Solutions. arXiv:2012.03165, 2020. \url{https://arxiv.org/abs/2012.03165}
\bibitem{ref55} Privacy Preservation Among Honest-but-Curious Edge Nodes: A Survey. arXiv:2208.05922, 2022. \url{https://arxiv.org/abs/2208.05922}
\bibitem{ref56} Quantifying Surveillance in the Networked Age: Node-based Intrusions and Group Privacy. arXiv:1803.09007, 2018. \url{https://arxiv.org/abs/1803.09007}
\bibitem{ref57} How BlockChain Can Help Enhance The Security And Privacy in Edge Computing?. arXiv:2111.00416, 2021. \url{https://arxiv.org/abs/2111.00416}
\bibitem{ref58} Private Edge Computing for Linear Inference Based on Secret Sharing. arXiv:2005.08593, 2020. \url{https://arxiv.org/abs/2005.08593}
\bibitem{ref59} SHEATH: Defending Horizontal Collaboration for Distributed CNNs against Adversarial Noise. arXiv:2409.17279, 2024. \url{https://arxiv.org/abs/2409.17279}
\bibitem{ref60} A Survey on Privacy-Preserving Caching at Network Edge: Classification, Solutions, and Challenges. arXiv:2405.01844, 2024. \url{https://arxiv.org/abs/2405.01844}
\bibitem{ref61} Resource Management in Fog/Edge Computing. arXiv:1810.00305, 2018. \url{https://arxiv.org/abs/1810.00305}
\bibitem{ref62} Joint Task Offloading and Resource Allocation for Multi-Server Mobile-Edge Computing Networks. arXiv:1705.00704, 2017. \url{https://arxiv.org/abs/1705.00704}
\bibitem{ref63} Sustainable Edge Computing: Challenges and Future Directions. arXiv:2304.04450, 2023. \url{https://arxiv.org/abs/2304.04450}
\bibitem{ref64} Collaborative Resource Management and Workloads Scheduling in Cloud-Assisted Mobile Edge Computing across Timescales. arXiv:2405.20560, 2024. \url{https://arxiv.org/abs/2405.20560}
\bibitem{ref65} Dynamic Spectrum Sharing for Load Balancing in Multi-Cell Mobile Edge Computing. arXiv:1910.04841, 2019. \url{https://arxiv.org/abs/1910.04841}
\bibitem{ref66} Joint UPF and Edge Applications Placement and Routing in 5G \&amp; Beyond. arXiv:2407.07669, 2024. \url{https://arxiv.org/abs/2407.07669}
\bibitem{ref67} Mobility Aware Edge Computing Segmentation Towards Localized Orchestration. arXiv:2110.07808, 2021. \url{https://arxiv.org/abs/2110.07808}
\bibitem{ref68} Energy Minimization for Mobile Edge Computing Networks with Time-Sensitive Constraints. arXiv:2008.09504, 2020. \url{https://arxiv.org/abs/2008.09504}
\bibitem{ref69} Online Learning for Offloading and Autoscaling in Renewable-Powered Mobile Edge Computing. arXiv:1609.05087, 2016. \url{https://arxiv.org/abs/1609.05087}
\bibitem{ref70} Energy-Aware Offloading in Time-Sensitive Networks with Mobile Edge Computing. arXiv:2003.12719, 2020. \url{https://arxiv.org/abs/2003.12719}
\bibitem{ref71} ENORM: A Framework For Edge NOde Resource Management. arXiv:1709.04061, 2017. \url{https://arxiv.org/abs/1709.04061}
\bibitem{ref72} Two-Stage Distributionally Robust Edge Node Placement Under Endogenous Demand Uncertainty. arXiv:2401.08041, 2024. \url{https://arxiv.org/abs/2401.08041}
\bibitem{ref73} AI-based Fog and Edge Computing: A Systematic Review, Taxonomy and Future Directions. arXiv:2212.04645, 2022. \url{https://arxiv.org/abs/2212.04645}
\bibitem{ref74} Retrieval-Augmented Generation for Mobile Edge Computing via Large Language Model. arXiv:2412.20820, 2024. \url{https://arxiv.org/abs/2412.20820}
\bibitem{ref75} 6GSoft: Software for Edge-to-Cloud Continuum. arXiv:2407.05963, 2024. \url{https://arxiv.org/abs/2407.05963}
\bibitem{ref76} 5G: Agent for Further Digital Disruptive Transformations. arXiv:1909.10152, 2019. \url{https://arxiv.org/abs/1909.10152}
\bibitem{ref77} Decentralized Blockchain-based model for Edge Computing. arXiv:2106.15050, 2021. \url{https://arxiv.org/abs/2106.15050}
\bibitem{ref78} Edge computing for cyber-physical systems: A systematic mapping study emphasizing trustworthiness. arXiv:2112.00619, 2021. \url{https://arxiv.org/abs/2112.00619}
\bibitem{ref79} Microservice-based edge platform for AI services. arXiv:2412.01328, 2024. \url{https://arxiv.org/abs/2412.01328}
\bibitem{ref80} Edge Learning for 6G-enabled Internet of Things: A Comprehensive Survey of Vulnerabilities, Datasets, and Defenses. arXiv:2306.10309, 2023. \url{https://arxiv.org/abs/2306.10309}
\bibitem{ref81} Edge Intelligence: Paving the Last Mile of Artificial Intelligence with Edge Computing. arXiv:1905.10083, 2019. \url{https://arxiv.org/abs/1905.10083}
\bibitem{ref82} Edge-Cloud Collaborative Computing on Distributed Intelligence and Model Optimization: A Survey. arXiv:2505.01821, 2025. \url{https://arxiv.org/abs/2505.01821}
\bibitem{ref83} Dissecting Open Edge Computing Platforms: Ecosystem, Usage, and Security Risks. arXiv:2404.09681, 2024. \url{https://arxiv.org/abs/2404.09681}
\end{thebibliography}

\end{document}
